% !TeX root = ../main.tex
\section{Limitations and Open Problems}
\label{sec:limitations}

\textbf{8.1 Instrument dependence.} All distributional gates are relative to
the fixed $E,D$; guarantees are for the scores as defined (Type-I end-to-end
by Prop.~\ref{prop:composite}; Type-II bounded by Prop.~\ref{prop:power} and
audited). Encoder-class invariance is open. The field study adds a concrete
instance: V4's hedged-reconciliation capability is contingent on the
instrument's hedge false-positive rate: zero for a calibrated
DeBERTa-large-MNLI (capability vacuous) and $2.3\%$ for a RoBERTa-large-MNLI
(capability active, $+0.035$ over \textsc{Wsum})---an instrument-conditional
dependence the replication demonstrates rather than merely hypothesizes.

\textbf{8.2 Assumption ledger.} Regime I requires exchangeability; Regime II
requires estimable shift; Regime III guarantees long-run frequency only.
Per-instance guarantees under fully adversarial shift are impossible in
general; the fail-safe posture (Prop.~\ref{prop:safeside}) is the stated
conduct in that gap. Marginal (not conditional) coverage is delivered;
exact conditional coverage is in fact unattainable distribution-free
\cite{barber2021}, so the designated next theoretical step is the achievable
relaxation: coverage conditional on prompt-difficulty strata via Mondrian
conformal. The field E2 audit adds: the
coverage-distribution refinement is not directly auditable by a naive KS of
empirical coverages (test-binomial noise, pool reuse, and score ties each
break the comparison---null-checked in \S\ref{sec:field-e2}); the exceedance
bound is tie-safe and held in all three runs.

\textbf{8.3 Relaxation gap and certificate-layer fragility.} A continuous
zero of $f_3$ certifies only relaxed consistency (Prop.~\ref{prop:gap}); the
certifier closes the gap at worst-case exponential cost, and
Prop.~\ref{prop:no-poly} shows this is unavoidable unless
$\mathrm P=\mathrm{NP}$. The field study adds a caution on the certificate
layer's \emph{input}: hard clauses assembled by thresholding NLI edges at
$w\ge 0.5$ inherit the instrument's false-positive edges, driving the
thresholded SAT arm to chance on real chains (\S\ref{sec:field}); certificates
therefore certify the thresholded graph, not the text, and a calibrated
edge-inclusion threshold (or weighted-clause certificates) is the registered
repair before the certificate layer's field role can carry weight.

\textbf{8.4 Scope of the field evidence.} \S\ref{sec:field} uses
template-inserted violations on a single corpus (GSM8K), now under two
architecture-distinct instruments; the localization claim is scoped to
compositional violations, and the V1/V2 regime is provably
counting-equivalent (Prop.~\ref{prop:struct-equiv}). An organic-error corpus
and MATH remain registered next steps (\S\ref{sec:protocol}(8)).

\textbf{8.5 Lexicographic repair order.} A convention; no optimality claim is
made.

\begin{center}
\fbox{\parbox{0.92\linewidth}{\textbf{Open Problem}
\refstepcounter{theorem}\label{open:tarski}
(Tarski unification). Recast the hard and soft logic layers as sheaves valued
in complete lattices with the Tarski Laplacian \cite{ghrist2022}, where global
sections are fixed points and Boolean semantics is native.}}
\end{center}

\medskip
\noindent\textbf{Open Problem (holonomy power theorem).} In a random-chain
model with a planted closing-claim violation of gap $\gamma$ and extraction
noise $\sigma$, derive the explicit power curve of the cycle-holonomy
statistic as a function of $\gamma/\sigma$ and the cycle's minimum edge
confidence, upgrading the observed field separation
(Table~\ref{tab:field}) into a predictive theorem.
